\section{Hanzo Models (HMMs)}
\label{sec:hmm}

\subsection{What Lives On Chain}

A Hanzo Model is a content-addressed, versioned bundle of weights,
metadata, and license. The chain does \emph{not} store weight blobs;
those live in content-addressed storage (Hanzo object storage,
\texttt{hanzo/storage}, or any compatible CAS) and are referenced by
their root hash. The chain stores a small fixed-size record per model
version that commits to everything a caller, an operator, or an auditor
needs to verify a job:

\begin{table}[h]
\centering
\small
\begin{tabular}{lp{6.0cm}}
\toprule
\textbf{Field} & \textbf{Definition} \\
\midrule
\texttt{model\_root}      & Merkle root over the weight tensors at
                            version $V$. Content-addressed; pinning the
                            root pins the weights bit-exactly. \\
\texttt{benchmark\_root}  & Merkle root over a canonical evaluation
                            harness output (e.g.\ MMLU, HELM, internal
                            evals). Submitted by the model owner;
                            challengeable on chain. \\
\texttt{license\_root}    & Hash of the license document (commercial,
                            research, or open under explicit terms),
                            with a machine-readable summary. \\
\texttt{inference\_price} & A price function $p(\cdot)$ in AI per
                            output token (or per call, per byte for
                            non-text models). May be a constant, a
                            piecewise function, or a reference to an
                            on-chain pricing oracle. \\
\texttt{training\_history}& Linked list of parent model\_roots, each
                            edge annotated with the training dataset
                            commitment, the trainer identity, and the
                            block height at which training completed. \\
\texttt{owner}            & Address authorized to publish new versions
                            and to update license / price; may be a
                            multisig or a DAO. \\
\bottomrule
\end{tabular}
\caption{HMM registry record. Six fields per model version. The chain
holds the commitments; weights and bench artifacts live in CAS.}
\label{tab:hmm-record}
\end{table}

\subsection{Versioning and Provenance}

Each new version of a model is a new record with a fresh
\texttt{model\_root} and a \texttt{training\_history} edge pointing to
the parent. This produces a directed acyclic graph rooted at the genesis
of each model family. Anyone can walk the graph to see the full
provenance: the original pretraining run, every fine-tune, every
distillation, and the dataset commitment used at each step.

\begin{figure}[h]
\centering
\begin{tikzpicture}[
    node distance=0.55cm and 1.1cm,
    box/.style={rectangle, draw, minimum width=1.9cm, minimum height=0.5cm,
                font=\footnotesize, rounded corners=1pt},
    arrow/.style={->, thick, gray}
]
\node[box] (v0) {hmm:zen-1.0};
\node[box, right=of v0] (v1) {hmm:zen-1.1};
\node[box, right=of v1] (v2) {hmm:zen-2.0};
\node[box, below right=of v0] (ft) {hmm:zen-1.0-medic};
\node[box, below=of v2] (q) {hmm:zen-2.0-q4};
\draw[arrow] (v0) -- (v1) node[midway,above,font=\tiny]{ft +1.2T tok};
\draw[arrow] (v1) -- (v2) node[midway,above,font=\tiny]{full retrain};
\draw[arrow] (v0) -- (ft) node[midway,below,font=\tiny]{specialize};
\draw[arrow] (v2) -- (q) node[midway,right,font=\tiny]{quantize};
\end{tikzpicture}
\caption{Provenance DAG. Each edge is a chain transaction recording the
parent model\_root, the training dataset commitment, and the trainer
identity.}
\end{figure}

\subsection{License Enforcement}

The \texttt{license\_root} commits to a license document with three
machine-readable fields: \emph{usage} (commercial / research / open),
\emph{redistribution} (allowed / restricted / prohibited), and
\emph{derivative} (allowed / restricted / prohibited). Compute operators
are required to read and respect the license at job time; the
\texttt{INFER\_REQUEST} precompile (\S\ref{sec:arch}) refuses to schedule
a job if the caller has not posted a matching license attestation.

\subsection{Price Function}

\texttt{inference\_price} is a first-class on-chain object. The simplest
form is a constant \texttt{uint256} in AI per output token. More
expressive forms are supported:

\begin{itemize}[leftmargin=1.2em]
  \item \emph{Piecewise tiered}: a price ladder that lowers per-token
        cost above batch thresholds.
  \item \emph{Time-of-day}: a function over block timestamp, used to
        clear demand at peak.
  \item \emph{Oracle-backed}: a reference to an on-chain price oracle
        whose own update mechanism is visible on chain.
\end{itemize}

The model owner sets the function and may update it; updates are
rate-limited (one update per epoch by default) to prevent surprise
price moves between job submission and execution.

\subsection{Benchmark Commitments and Challenges}

A model owner publishes a \texttt{benchmark\_root} when registering a
version. Any address may post a \emph{challenge}: a deposit-bonded
claim that the published benchmark output does not match the actual
output of the published \texttt{model\_root} on the canonical eval
harness. Challenges are resolved by a designated set of operators
re-running the harness inside a TEE and signing the result. If the
challenge succeeds, the challenger collects the model owner's
benchmark bond; if it fails, the challenger forfeits theirs. This
matches the dispute pattern used elsewhere in Lux for attestation
challenges (LP-134 §6).

\subsection{Native Hanzo Models at Activation}

At 25 December 2025 activation, the chain is genesis-seeded with the
core Hanzo model family registered at version 1.0:

\begin{itemize}[leftmargin=1.2em]
  \item \texttt{hmm:zen-text} (text-to-text foundation model);
  \item \texttt{hmm:zen-vision} (multimodal vision-language);
  \item \texttt{hmm:zen-audio} (speech and audio);
  \item \texttt{hmm:zen-code} (code-specialized);
  \item \texttt{hmm:zen-embed} (embedding model).
\end{itemize}

These are the production models served by the Hanzo LLM Gateway
(\texttt{\textasciitilde/work/hanzo/llm}) and by Jin
(\texttt{\textasciitilde/work/hanzo/jin}). Registration on chain is the
single source of truth for which weights are canonical and at what price
they are offered.
